\begin{frame}{똑똑한 초기화: k-means++}
  무작위 초기화의 결함 구조: $k$개 중심점을 데이터에서
  \textbf{독립적으로} 무작위로 뽑으면, 여러 중심점이
  \textbf{같은 좁은 영역에 몰리는} 가능성을 스스로 만듦.
  k-means++는 이 결함을 정면으로 고친다:
  \vskip0.4em
  \begin{enumerate}\small
    \item 첫 중심점: 데이터에서 \textbf{균등 무작위}
    \item 각 점 $x$에 대해 이미 고른 중심점까지의 \textbf{최소
      거리} $D(x)$ 계산 $\to$ \textbf{$D(x)^2$에 비례하는 확률}로
      다음 중심점 선택
    \item $k$개 다 고를 때까지 2 반복
  \end{enumerate}
  \vskip0.6em
  핵심: ``가장 가까운 중심점에서 \textbf{멀리} 있는 점을 더 높은
  확률로 고른다'' --- 첫 중심점이 왼쪽 덩어리에 놓이면 두 번째는
  오른쪽 덩어리에서 상대적으로 높은 확률로 추첨 $\to$
  \textbf{같은 덩어리에 두 중심점이 몰리는 상황을 구조적으로
  막음}.
  \vskip0.4em
  왜 \textbf{제곱}? Arthur \& Vassilvitskii(2007)의
  k-means++ 논문: $D(x)^2$-비례 선택을 쓸 때, 최종 분할의 SSE가
  \textbf{기댓값 기준으로} 최적 분할 SSE의 $O(\log k)$배
  이내(원논문 bound $8\ln k + 7$배)에 머문다는
  \textbf{증명 가능한 근사 보장} --- 단순 무작위에는 없는 성질.
\end{frame}
